\subsection{The actual periodic theta embedding and its marked coordinates}
\label{subsec:ftb-actual-theta-embedding}

We retain the source's large label set.  Put
\[
 I_M=(\mathbb Z/M\mathbb Z)^{2g},\qquad
 [a,b]=\sum_{i,j}a_ib_jJ_{ij}\in\mathbb Z/M\mathbb Z,
 \tag{TE1}\label{eq:ftb-theta-labels}
\]
where \(J\) is the displayed standard alternating matrix.
An actual marked object consists of a principally polarized
abelian scheme \(A/S\), its symmetric rigidified ample
principal line bundle \(L\), endomorphisms if prescribed,
and a full level-\(M\) basis \(p_i\).
Thus \(\ell(a)=\sum_i a_ip_i\) identifies geometric
\(M\)-torsion with \(I_M\), and its Weil pairing is
\[
                         e_M(\ell(a),\ell(b))=\zeta^{[a,b]}.
 \tag{TE2}\label{eq:ftb-labelled-weil}
\]
The base contains the specified primitive root \(\zeta\),
and \(M\) is invertible there.

Let \(N'\) be odd, \(m\mid N'\), \(r=N'/m\), and
\(M=N'^2=m^2r^2\).
The source datum \(\operatorname{Emb}(m)\) is the
following actual section system, not a choice of
projective coordinates made after base change.
Choose a generator \(\theta\) of the rank-one module
of sections of \(L\), the cube/symmetry multiplication
isomorphism \([m]^*L\simeq L^{m^2}\), and its resulting
section \(s_0\) of \(L^{m^2}\).
For every \(a\in I_M\), take the odd-order symmetric
theta-group lift \(U_a\) of
\(\ell(a)^{r^2}\in A[m^2]\), and put
\[
                                \sigma_a=U_a s_0.
 \tag{TE3}\label{eq:ftb-embedding-sections}
\]
Here group powers denote multiplication on the abelian
scheme, and the theta-group action includes its actual
line-bundle isomorphism.

The abelian-scheme laws used in constructing these data
have the same scopes as in the key construction:
the module of sections of \(L\) is locally free of rank
one and its displayed local bases commute with arbitrary
base change; the cube and symmetry isomorphisms give
\([m]^*L\simeq L^{m^2}\); odd invertible torsion has
locally existing, unique symmetric lifts; and on a
geometric field fibre \(h^0(L^q)=q^g\), \(L^q\) is very
ample for \(q\ge4\), and the Weil pairing on \(q\)-torsion
is perfect.
These are used on the same \(A,L\).

\begin{lemma}[Existence, spanning and periodicity]
\label{lem:ftb-actual-periodic-embedding}
For every point of an affine test base, the data
\eqref{eq:ftb-embedding-sections} exist on an affine open
neighbourhood of that point.
They commute, with their section-coordinate vectors,
with maps of affine bases.
If \(m\ge2\), on every algebraically closed field fibre
they span all sections of \(L^{m^2}\), give its actual
projective coordinate representation, and satisfy
\[
             a\equiv a'\pmod {m^2}
                         \Longrightarrow \sigma_a=\sigma_{a'}.
 \tag{TE4}\label{eq:ftb-theta-section-periodicity}
\]
Every \(N'\)-torsion point has the form \(\ell(h)^{N'}\).
On such a fibre the coordinate line is exactly \(L^{m^2}\);
in particular
\[
                           h^0((L^{m^2})^e)=(m^2)^g e^g.
 \tag{TE5}\label{eq:ftb-theta-exact-h0}
\]
\end{lemma}
\begin{proof}
Choose a neighbourhood on which the rank-one section
module has a generator.
For each of the finitely many \(a\), choose a
neighbourhood carrying its symmetric lift, and
intersect them.
A basic affine neighbourhood inside this intersection
gives the first assertion.
The multiplication isomorphism pulls to this base;
pulling \(\theta\) through it and applying the finitely
many lifts gives the sections.
This is literally the source's common-neighbourhood
construction.

We spell out the nontrivial base-change assertion for
the generator.
For affine bases \(T'\to T\), let \(\theta\) generate
\(H^0(A_T,L)\).
Locally on \(T\), choose a base-change-compatible
rank-one basis \(b\).
The proper scheme \(A_T\) is quasi-compact and
quasi-separated.
A section on a principal base open extends after
multiplication by a power of its defining function;
a section vanishing there is killed by another power.
Indeed, express it on finitely many affine line charts,
clear their finitely many denominators, and clear the
finitely many overlap differences.
This is the localization theorem for the quasi-coherent
invertible sheaf on this actual proper scheme.
It lets one write a power of the local \(b\) as a scalar
multiple of \(\theta\).
Pull back the compatible basis to \(T'\), express any
given section in it, clear its scalar denominator,
and clear the resulting section difference.
Near each point of \(T'\) this gives a scalar \(c\),
invertible at that point, such that \(c\) times the
section lies in the span of the pulled-back \(\theta\).
The ideal of all such \(c\) is the unit ideal.
A finite linear combination of them equal to \(1\)
shows that the pulled-back \(\theta\) generates globally.
The multiplication and lift isomorphisms themselves
pull back by their explicit line-chart products.
Consequently \(\sigma'_a\) is the actual pulled-back
\(\sigma_a\); their local germs agree.
Pulling a line generator at a test point proves the
same section-coordinate condition in both directions.
This is stronger than comparison just at reduced points.

For spanning, use the source's finite Heisenberg
linear-algebra proof.
Let \(V=H^0(A,L^q)\) over an algebraically closed field,
\(q=m^2\), and choose a lift \(U_x\) of each \(x\in A[q]\).
The operators are invertible and
\[
                         U_xU_y=e_q(x,y)U_yU_x .
 \tag{TE6}\label{eq:ftb-heisenberg-operators}
\]
They are linearly independent in \(\operatorname{End}V\).
To verify this, take a dependence with minimal finite
support and a nonzero coefficient at \(x_0\).
Commute it with \(U_y\) and subtract its multiple
\(e_q(x_0,y)\).
The coefficient at \(x_0\) disappears.
Minimality implies that every remaining nonzero
coefficient has
\(e_q(x,y)=e_q(x_0,y)\) for every \(y\).
Perfectness gives \(x=x_0\).
The resulting one-term dependence contradicts
invertibility of \(U_{x_0}\).
Since
\[
            \#A[q]=q^{2g}=(\dim V)^2=\dim\operatorname{End}V,
\]
these operators span every endomorphism.
Any nonzero invariant subspace is therefore all of \(V\):
given \(0\ne v\) in it and any \(w\in V\), apply the
rank-one operator \(u\mapsto\varphi(u)w/\varphi(v)\),
where \(\varphi(v)\ne0\).

The section \(s_0\) is nonzero.
Otherwise surjectivity of \([m]\), and integrality of
the geometric abelian fibre, would make \(\theta=0\),
contrary to \(h^0(L)=1\) and its generating property.
The points \(\ell(a)^{r^2}\) cover \(A[m^2]\).
Indeed, in the free group \(I_M\), an element killed
by \(m^2\) has every coordinate divisible by \(r^2\);
divide those coordinates and use the full marked basis.
Thus the span of the \(\sigma_a\) is invariant under
all theta-group operators and is nonzero.
It equals \(V\).

Since \(m^2\ge4\), the complete section system of
\(L^{m^2}\) embeds the geometric abelian variety.
Every one of those complete sections is a linear
combination of the \(\sigma_a\).
The same linear combinations show that this specified
system gives the embedding.
Its actual projective-coordinate representation and
the affine-ratio very-ampleness condition follow by
pulling the standard coordinate charts of projective
space.
This argument preserves the coordinate line, as well
as the points.

For periodicity write \(N'=mr\).
If \(a-a'\) is divisible by \(m^2\), then
\(r^2(a-a')=0\) in \(I_M\).
The two lifted translation points in
\eqref{eq:ftb-embedding-sections} are equal.
Uniqueness of their odd-order symmetric lifts gives
the same action on \(s_0\), proving
\eqref{eq:ftb-theta-section-periodicity}.
Finally, \([N']\) maps \(A[N'^2]\) onto \(A[N']\);
equivalently, divide the coordinates of an element
killed by \(N'\) in the full marked basis.
This gives the last torsion assertion.
The coordinate-line isomorphism and the principal
section-rank formula give
\eqref{eq:ftb-theta-exact-h0}.
\end{proof}

\paragraph{Projective representation over the local base.}
The same locally constructed system is projective,
with its all-ring-test coordinate property, before
specializing to a geometric field.
Here is the source's argument for this upgrade.
The geometric spanning result implies that its
coordinates have no common zero on any geometric
fibre, hence no common zero on the whole proper scheme.
They give a morphism \(f:A_T\to\mathbb P_T^{I_M-1}\)
whose coordinate-line comparison is the actual one.
It is proper.
Two geometric points with the same image lie in the
same geometric fibre and are equal by its coordinate
representation.
Likewise a dual-number point whose image is constant
has constant base point; it lies in that geometric
fibre and is constant by its all-test-ring coordinate
uniqueness.
The tangent criterion makes \(f\) formally unramified.
Its diagonal is then an open immersion.
Geometric injectivity makes that diagonal universally
surjective, hence an isomorphism.
Thus \(f\) is a monomorphism, and properness makes it
a closed immersion.
The coordinate representation of this closed
subscheme of projective space is exactly the
all-test-ring \(\operatorname{ProjRep}\) for its
own ideal.

We next give the coordinate formula which allows this
same embedding to be reconstructed from a marked theta
vector.  No independently normalized theta function
is used.

\begin{proposition}[The source's torsion coordinates]
\label{prop:ftb-actual-torsion-theta-coordinates}
Assume \(m^2\mid N'\), \(m\ge2\), and let \(v=(v_a)_{a\in I_M}\)
be an actual \(\operatorname{HasTh}\) vector for the
marked \(A,L\) over an affine base \(R\).
Assume \((v_a)_a=R\).
For the system \eqref{eq:ftb-embedding-sections}, put
\[
 \begin{split}
 \operatorname{idx}(a,h)&=ra+mh\in I_M,\\
 E(a,h)&=\left(\frac{M+1}{2}[mh,ra]\right)_{\!M}
                                                    \in\{0,\ldots,M-1\},\\
                         p_h(a)&=\zeta^{E(a,h)}v_{\operatorname{idx}(a,h)} .
 \end{split}
 \tag{TE7}\label{eq:ftb-actual-torsion-coordinate-formula}
\]
Here the subscript \(M\) takes the least residue; the
integer \((M+1)/2\) is the inverse of \(2\) modulo the
odd number \(M\).
On an affine open neighbourhood of every point of
\(R\), \(p_h\) is an actual section-coordinate vector
of \((\sigma_a)\) at \(\ell(h)^{N'}\).
This holds over arbitrary test rings, including rings
with nilpotents, and is compatible with affine
base change.
\end{proposition}
\begin{proof}
We expand the scalar calculation determining
\eqref{eq:ftb-actual-torsion-coordinate-formula}.
If \(q=2l+1\) and symmetric \(q\)-order lifts of \(x,y,x+y\)
act in a fixed line chart, write their unit functions
as \(f,g,e\).
Their product differs from the lift of \(x+y\) by one
base unit \(d\).
Write \(c=e_q(x,y)\), so
\[
                   f(z)g(x+z)=c\,g(z)f(y+z)
                              =d\,e(z).
\]
The \(q\)-fold product identity for each symmetric
lift, and iteration of the first equality, give
\(c^q=1\) and \(d^q=1\).
Their symmetry identities, evaluated at
\(0,-x,-x-y\), give \(d^2=c\):
substitute these three identities in the two displayed
product formulas and cancel their invertible factors.
Consequently
\[
                         d=c^{l+1}=c^{(q+1)/2}.
 \tag{TE8}\label{eq:ftb-symmetric-lift-half-commutator}
\]
This is the source's \(\mathrm{scalar\_eq}\) calculation.
It takes place in groups of units and never divides
by \(q\), by \(2\), or by a theta value.

The multiplication isomorphism between levels pulls a
lift through \([r]\) by conjugation.
In local unit functions, with its comparison generator
\(\beta\), the resulting lift has function
\[
                   u'(z)=\beta(z)\,u([r]z)\,
                                         \beta(y+z)^{-1}.
 \tag{TE9}\label{eq:ftb-lift-pullback-conjugation}
\]
Its \(q'\)-fold product is one: the \(\beta\) factors
telescope, the old \(q\)-fold product is one, and
\(q\mid q'\).
Its symmetry condition follows by the same cancellation.
Thus this is the actual symmetric lift at the new
level, not just an equality of its translation points.

Take \(q'=M\), \(Y=\ell(ra)\), \(Z=\ell(mh)\).
Then \([r]Z=\ell(h)^{N'}\), and the Weil scalar in
\eqref{eq:ftb-symmetric-lift-half-commutator} is
\(\zeta^{[mh,ra]}\) by
\eqref{eq:ftb-labelled-weil}.
Apply \eqref{eq:ftb-lift-pullback-conjugation}
to the two multiplication-pullback theta sections
occurring in the definition of \(v\).
The acted section at \(Z\), followed by evaluation at
the origin, is the section labelled \(ra+mh\).
The line-chart formulas for action, multiplication
pullback and evaluation are respectively
\[
 (U_xs)(z)=u_x(z)s(x+z),\qquad
              [r]^*s(z)=\beta(z)s([r]z),\qquad
                             s(z)=\tau_z\,\operatorname{value}_z(s).
\]
Substituting these three formulas and
\eqref{eq:ftb-symmetric-lift-half-commutator} cancels
all chart factors and gives one unit \(u_h\), independent
of \(a\), with
\[
          \operatorname{value}_{\ell(h)^{N'}}(\sigma_a)
                           =u_h\,\zeta^{E(a,h)}v_{ra+mh}.
 \tag{TE10}\label{eq:ftb-common-unit-torsion-values}
\]
The independence is exact: its remaining factors are
the multiplication comparison generators and the
chosen generator of \(L^{m^2}\) at \(\ell(h)^{N'}\);
none involves \(a\).

We also check that the theta section chosen for the
embedding can be compared by a unit to the section in
the given \(\operatorname{HasTh}\) datum.
Both generate the same rank-one section module after
affine base change.
If \(\theta=\rho\theta'\) and
\(\theta'=\rho'\theta\), then
\((1-\rho\rho')\theta=0\).
Action, multiplication pullback and evaluation show
that any scalar annihilating \(\theta\) annihilates
every \(v_a\).
Their unit-ideal hypothesis makes that scalar zero.
Thus \(\rho\rho'=1\), and \(\rho\) is a unit.
Its factor in \eqref{eq:ftb-common-unit-torsion-values}
can be absorbed into the same common \(u_h\).
This is why the proposition retains the full
unit-ideal hypothesis.

Choose an affine neighbourhood where the original
theta-point data and the line generator at
\(\ell(h)^{N'}\) exist.
Equation \eqref{eq:ftb-common-unit-torsion-values}
then says that the pulled-back sections equal
\(u_h p_h(a)\) times that line generator.
Replace the generator by \(u_h\) times itself.
This is precisely the section-coordinate condition
for \(p_h\).
All operations above are operations on actual
invertible line-chart functions and their pullbacks.
They prove the assertion on arbitrary rings.
The affine base-change comparison in the preceding
lemma preserves these equalities.
\end{proof}

\paragraph{The bounded torsion-form application.}
Take an odd prime \(n\), \(N'=n^j\), \(m=n^2\), and \(j\ge4\).
Then \(m^2=n^4\mid N'\).
The small coordinate set for
Proposition~\ref{prop:ftb-uniform-torsion-forms} is
\[
                  I_{n^4}=(\mathbb Z/n^4\mathbb Z)^{2g},
\]
and the large set is \(I_{n^{2j}}\), with their
coordinatewise reduction map.
By \eqref{eq:ftb-theta-section-periodicity} the
sections repeat through that reduction.
Their actual line bundle is \(L^{n^4}\);
\eqref{eq:ftb-theta-exact-h0} gives
\(\delta=(n^4)^g\), so use \(C=n^4\).
The marked family is
\(\ell(h)^{n^j}\), \(h\in I_{n^{2j}}\), covering every
\(n^j\)-torsion point.
Therefore one \(j_0\), depending only on \(g,n\)
(and harmless fixed marking data), supplies \(d,Q\)
and the source's \(\operatorname{CutBy}\) conclusion
for all these actual geometric theta systems.
For instance choose \(j_0\ge4\) and at least the
\(k_0\) in \eqref{eq:ftb-torsion-cut-choices};
\(n^j>j\) then gives \(n^j\ge k_0\).
This is the actual instantiation of the previously
supplied-system degree bound.

\paragraph{Source concordance.}
The finite Heisenberg proof is
\texttt{AbelianSchemeLawProofs/LinearIndependentOfCommSmul.lean}
(\texttt{ff0a2b43b43c5f2e52dbea36aaa1a5183ab03854});
its actual Weil/rank application is
\texttt{IsFieldGlobalFnsSpec.lean}
(\texttt{0a7262ebe77a38822f8a01674c74d0321f157569}).
The local construction, base change, geometric spanning
and projective representation are
\texttt{ModuliProofs/ExistsIsGenSecNear.lean}
(\texttt{b4aa72654a37fcd1c350867471395b8c0d010d6a}),
\texttt{ThetaEmbeddingSectionsBaseChangeAffine.lean}
(\texttt{2731c3145c87abc0803904f135aab2919d67aa36}),
\texttt{EmbeddingSectionsSpanGeometricFibre.lean}
(\texttt{a049aa002403cae28d0d21e1784f9881e9ae60d2}),
\texttt{ProjCoordRepresentationGeometricFibreSpanning.lean}
(\texttt{8959fedebf0c889d3d6f2a438ecf1f42ebf844d3}),
and \texttt{MonoOfGeomPoints.lean}
(\texttt{dbc88a67769f1e74f53d167c4167d098c69470d0}).
Periodicity is
\texttt{ThetaSectionsVeryAmplePeriodic.lean}
(\texttt{396af0ba916b99e5242f0b2d8e81298c5f5ebe03}).
The scalar/conjugation formulas and marked coordinates
are
\texttt{ThetaFamilyOrigin/IsThetaPtOverOfSubsingleton.lean}
(\texttt{915c8b61564aa13483d4e060433e1a248b2d11fa}),
\texttt{IsVVGenOfIsThetaFamily.lean}
(\texttt{15448cd6bfcaa6cea793423b43c327b42ca17631}),
and \texttt{ModuliProofs/TorsionSectionCoordsThetaTimesRoots.lean}
(\texttt{4d8fa964670d2fe0a868543c17f1f5e699e58977}).
The final degree application is
\texttt{EmbeddedAbelianVarietyCutOutByTorsionForms.lean}
(\texttt{49df7ae02a38038f9df4ce1fa30697e649ef5062}).
The common path prefix is
\texttt{src/Development/AbelianVarietyTheory/Proofs/}.
Addition and endomorphism graph systems, and the
subsequent very-dense-family reconstruction, remain
separate next producers.
